\documentclass[12pt]{article}

\usepackage[margin=1cm, paperwidth=7cm, paperheight=12cm]{geometry}
\usepackage{fontspec}
\usepackage{polyglossia}
\usepackage{luabidi}
\usepackage{xcolor}

\setmainlanguage{arabic}
\setotherlanguage{english}
\newfontfamily\arabicfont[Script=Arabic, Scale=1.0]{Amiri}
\newfontfamily\englishfont{Latin Modern Roman}
\newcommand{\al}{\textarabic}
\newcommand{\el}{\textenglish}

\newcommand{\Hl}[1]{\colorbox{yellow!40}{\strut #1}}

\pagestyle{empty}

\title{
  موقف سيدنا عمر -رضي الله عنه- يوم السقيفة
  }
\date{}
\begin{document}
% \maketitle
\centering
\noindent\textbf{موقف سيدنا عمر بن الخطاب -رضي الله عنه- يوم السقيفة}\\
\vspace{20pt}
  \al{
  اجتمع الصحابة في سقيفة بني ساعد لختيار خليفة للمسلمين، فسمع عمر -رضي الله عنه- جلبة في الباب؛ فقيل له: بعض الأعراب جاؤوا يشاركون في الأمر. فخرج لهم الفاروق رضي الله عنه وقال:\strut
  }
  \al{\Hl{\strut
  ليرجع صاحب المحراث إلى
  }}
  \al{\Hl{\strut
   محراثه وصاحب الصنعة إلى صنعته، 
  }}
  \al{\Hl{\strut
   الأمر اليوم للمهاجرين والأنصار.
  }}
\end{document}